\chapter{Germination, Definition and Description}

In the rate studies described herein, germination was taken as the time when
the radicle (or rarely the epicotyl) could first be detected breaking through the seed
coat. In general the radicle elongates rapidly so that there was little ambiguity in
deciding when a seed had germinated. Less common were confusing situations
where the seed would expand 10-40\% splitting the seed coat, but the radicle would
not start development until after a significant period of time or after a shift to another
temperature. In these situations germination was not recorded until the radicle began
to develop. This latter behavior was found in Comus, Corydalis, Daphne, Euonymus,
Prunus, and Taxus.

An important distinction is now proposed. \textit{One-step germinators} form the
photosynthesizing leaves (cotyledons or true leaves) within a few weeks of
germination and at the same temperature at which germination occurred. \textit{Two-step
germinators} form the radicle first. The radicle may or may not develop a bulb or corm
with attached root. The photosynthesizing leaves do not form until (a) after a
significant time gap of two or more months or (b) after one or two shifts of temperature.
Examples of the first behavior are Abeliophyllum distichum and Clematis lanuginosa
where simply keeping the seedlings fo(two to three months at the original germination
temperature of 70 will initiate development of leaves. More often a shift in temperature
is needed. Typical is the behavior of Trillium and about half of Lilium. With these a
radicle and bulb or corm forms at 70. When this is complete the seedlings have to be
shifted to temperatures around 35-40 for three months. On returning to 70 the
photosynthesizing leaf (usually a single leaf) forms which is a cotyledon in Trillium and
a true leaf in Lilium. The three months at 35-40 are a period where chemical changes
are occurring, and these temperatures are the conditions where \textit{these chemical
changes occur most rapidly}. The resumption of growth usually occurs on a shift to
warmer temperatures, but many examples were found where growth resumed at the
low temperature after a period of time and after the inhibitors had been destroyed.

These chilling periods are termed conditioning or conditioning periods.
Traditionally they have been described as dormancy or breaking dormancy. These
traditional terms imply inactivity when in fact certain processes are occurring at their
\textit{maximum metabolic rate}. The chilling period has also been described as stratification,
an archaic term derived from the old English .process of placing seeds in layers in
sand and placing them outdoors overwinter. This term is equally misleading as are
terms like vernalization or winterization. It is best to get to the heart of the matter and
call these periods \textit{conditioning periods} in the sense that they get the seeds into a
condition where they can germinate.

Flowering plants are divided into the monocotyledons with one germinating leaf
and dicotyledons with a pair of germinating leaves. Even this most useful of
divisions has an occasional exception. In the present work Claytonia virginica,
a dicotyledon, developed only a single germinating leaf. This was discussed
with Prof. Carl Keener at Penn State who said that Ranunculus ficaria also does
this and that the phenomenon is known. Orchidaceae never develop any
cotyledons. A description of their fascinating life cycle can be found in
Summerhayes (ref. 8).

Flowering, plants are also divided into epigeal and hypogeal germinators. In
epigeal germination the cotyledon(s) are the first photosynthesizing leaves. In
hypogeal germination the cotyledons wither within the seed coat, and the first first
photosynthesizing leaf or leaves are true leaves. Both monocotyledons and
dicotyledons have species with both types of germination. Both types are found within
the same family and even the same genus as in Lilium. In beans and peas the
cotyledons do some photosynthesis, but are primarily storage organs. Thus the
division into epigeal and hypogeal is not always clear and in any event does not seem
fundamental to plant classification. \textit{Most important}, epigeal and hypogeal must not be
confused with one-step and two-step germinators. One-step germinators can be
epigeal (most dicotyledons) or hypogeal (Iridaceae and Poaceae). Two-step
germinators can likewise be either epigeal (Trillium) or hypogeal (many Lilium). All
four combinations can be found in both monocotyledons and dicotyledons.

Before leaving germination it is appropriate to discuss percent germination and
the difficulties that arise in determing this number. Percent germination is the number
that has been used traditionally to measure germination and the effectiveness of
germination procedures. It is defined as the number germinated x 100 divided by the
total number of seeds. The number germinated can be determined with reasonable
certainty using the definition of germination given above. It is the determination of the
total number of seeds that causes the problem as illustrated by the following example.
Anyone who has cracked open a number of pits of a stone fruit like plums will note that
a number of the pits are hollow shells. Judging from external appearance these would
have been counted as seeds. Suppose the investigation is more sophisticated and
examines the seeds by X-rays and determines the number of such empty shells. It
does not stop there. Techniques could be developed that would reveal that some of
the seeds lacked embryos or had abnormal embryos. Thus the total number of seeds
will depend on the technique used to determine the number of seeds.

For these reasons percent germination are best used in comparing sequential
germination in a single sample or comparing results from different treatments of a
single sample. It is of much less significance in an absolute sense.

